%!TEX root = ../main.tex

\chapter{Einleitung}
Die Einleitung soll den Ausgangspunkt der Arbeit umreißen, in kurzer Form zur Problemstellung hinführen und das Interesse der lesenden Person für die Arbeit wecken. Allgemeine Einleitung ins Thema, keine Unternehmens- oder Produktbeschreibungen, Organigramme u.ä., wenn diese nicht direkt zum Thema führen. Ziele und Vorgehensweise nicht vermischen.

\section{Ausgangssituation}
ATLAS ist eine Webanwendung im Unternehmenskontext. In diesem Abschnitt wird die aktuelle Ausgangslage beschrieben: Welche Funktionen bereits vorhanden sind, wie der bisherige Entwicklungsstand ist und wie weit die automatisierten Tests aktuell fortgeschritten sind.

\textbf{Leitfragen:}
\begin{itemize}
	\item Was ist ATLAS in ein bis zwei Sätzen?
	\item Welche Testarten gibt es heute bereits?
	\item Wo bestehen aktuell Lücken in der Absicherung?
\end{itemize}

\section{Problemstellung und Zielsetzung}
Die Problemstellung ergibt sich aus dem Spannungsfeld zwischen wachsender Systemkomplexität und unzureichender Testabdeckung. Ziel ist es, eine belastbare Testinfrastruktur zu schaffen, die lokal nutzbar ist und perspektivisch zentral betrieben werden kann.

\textbf{Leitfragen:}
\begin{itemize}
	\item Warum reicht der bisherige Testansatz nicht aus?
	\item Welches konkrete Ziel soll mit dem Projekt erreicht werden?
	\item Welche Ergebnisse gelten als erfolgreich?
\end{itemize}

\section{Motivation}
Die Motivation liegt in einer höheren Entwicklungsqualität, schnelleren Feedbackzyklen und einer stabileren Zusammenarbeit im Team.

\subsection{Warum ist es wichtig?}
Mit zunehmendem Funktionsumfang steigt das Risiko für Regressionen. Ein systematischer Testansatz reduziert dieses Risiko und erhöht die Zuverlässigkeit von Releases.

\subsection{Welcher Mehrwert entsteht?}
Der Mehrwert zeigt sich in früher Fehlererkennung, besserer Wartbarkeit und einer objektiveren Bewertung der Softwarequalität über definierte Testläufe.

\subsection{Für wen entsteht der Mehrwert?}
Vom verbesserten Testprozess profitieren Entwicklerteam, Hauptentwickler, Stakeholder und Endnutzer durch stabilere Funktionalität und planbarere Auslieferungen.

\section{Aufbau der Arbeit}
Diese Arbeit gliedert sich in sechs Kapitel: Einleitung, Grundlagen, Ist-Analyse und Anforderungen, Konzeption der Testinfrastruktur, Umsetzung sowie kritische Reflexion mit Ausblick.
